\section{Bố cục của khóa luận}

Bố cục của khóa luận được trình bày thành 3 chương. Mỗi chương trình bày và thảo luận các vấn đề khác nhau liên quan đến khóa luận và được tóm tắt như sau:

\textbf{Chương 1: Giới thiệu.} 

Tổng quan về các vấn đề liên quan đến khóa luận, những thông tin về mạng cảm biến không dây và mạng ad-hoc, tìm hiểu về các thuật toán định tuyến phổ biến và ưu nhược điểm của từng thuật toán.

% \textbf{Chương 2: Webots và E-puck robot hai bánh vi sai} 

% Trình bày tổng quan về trình mô phỏng vật lý Webots  và cấu hình chi tiết của mô hình robot E-puck. Trình bày mô hình điều khiển động học hai bánh vi sai.

\textbf{Chương 2: Cơ sở lý thuyết} 
Trình bày về các khái niệm liên quan đến mạng cảm biến không dây và mạng ad-hoc, các giao thức định tuyến trong mạng ad-hoc và thuật toán định tuyến dựa trên Gradient. Bên cạnh đó cũng giới thiệu về nền tảng phần cứng nRF54L15 và giao thức Bluetooth Low Energy Mesh.

\textbf{Chương 3: Triển khai trên phần cứng nRF54L15} 

Trình bày về quá trình thiết kế và triển khai firmware cho vi xử lý nRF54L15, kiến trúc phần mềm tổng thể của hệ thống và các hàm chức năng chính được triển khai.

\textbf{Chương 4: Kết quả và thảo luận.} 

Trình bày các kết quả thí nghiệm khi triển khai trên mạng thực tế, so sánh và đánh giá hiệu suất với các giải pháp khác và thảo luận về những thách thức gặp phải trong quá trình phát triển.

\textbf{Kết luận và hướng phát triển.} 
Các tổng kết về các công việc đã thực hiện trong khóa luận và các kết quả đạt được. Đồng thời cung cấp các hướng phát triển trong tương lai dựa trên kết quả nghiên cứu của khóa luận.